%-----------------------------------------------------------------------------------------
% Autor dieser Vorlage:
% Stefan Macke (http://fachinformatiker-anwendungsentwicklung.net)
% Permalink zur Vorlage: http://fiae.link/LaTeXVorlageFIAE
%
% Sämtliche verwendeten Abbildungen, Tabellen und Listings stammen von Dirk Grashorn.
%
% Lizenz: Creative Commons 4.0 Namensnennung - Weitergabe unter gleichen Bedingungen
% -----------------------------------------------------------------------------------------

\documentclass[
	ngerman,
	toc=listof, % Abbildungsverzeichnis sowie Tabellenverzeichnis in das Inhaltsverzeichnis aufnehmen
	toc=bibliography, % Literaturverzeichnis in das Inhaltsverzeichnis aufnehmen
	footnotes=multiple, % Trennen von direkt aufeinander folgenden Fußnoten
	parskip=half, % vertikalen Abstand zwischen Absätzen verwenden anstatt horizontale Einrückung von Folgeabsätzen
	numbers=noendperiod % Den letzten Punkt nach einer Nummerierung entfernen (nach DIN 5008)
]{scrartcl}
\pdfminorversion=5 % erlaubt das Einfügen von pdf-Dateien bis Version 1.7, ohne eine Fehlermeldung zu werfen (keine Garantie für fehlerfreies Einbetten!)
\usepackage[utf8]{inputenc} % muss als erstes eingebunden werden, da Meta/Packages ggfs. Sonderzeichen enthalten

% !TEX root = Projektdokumentation.tex

% Hinweis: der Titel muss zum Inhalt des Projekts passen und den zentralen Inhalt des Projekts deutlich herausstellen

\newcommand{\CBP}{Computergestützte Betriebsprüfung}
\newcommand{\GA}{Globales Aktualisieren}

\newcommand{\titel}{\GA\space von Dokumenten}
\newcommand{\untertitel}{\CBP\space - Abschluss \& Dokumente}
\newcommand{\kompletterTitel}{\titel{} -- \untertitel}

\newcommand{\autorName}{Guido Eckelt}
\newcommand{\autorAnschrift}{Boddinstraße 30}
\newcommand{\autorOrt}{12053 Berlin}

\newcommand{\betriebLogo}{LogoDRVBund.png}
\newcommand{\betriebLogoMitTitel}{LogoDRVBundMitTitel.png}
\newcommand{\betriebName}{\textsc{Deutsche Rentenversicherung} Bund}
\newcommand{\betriebAnschrift}{Ruhrstraße 2}
\newcommand{\betriebOrt}{10704 Berlin}

\newcommand{\ausbildungsberuf}{Fachinformatiker für Anwendungsentwicklung}
\newcommand{\betreff}{Dokumentation zur betrieblichen Projektarbeit}
\newcommand{\pruefungstermin}{Winter 2016}
\newcommand{\abgabeOrt}{Berlin}
\newcommand{\abgabeTermin}{20.12.2016}
 % Metadaten zu diesem Dokument (Autor usw.)
% !TEX root = ../Projektdokumentation.tex

% Anpassung an Landessprache ---------------------------------------------------
\usepackage{babel}

% Umlaute ----------------------------------------------------------------------
%   Umlaute/Sonderzeichen wie äüöß direkt im Quelltext verwenden (CodePage).
%   Erlaubt automatische Trennung von Worten mit Umlauten.
% ------------------------------------------------------------------------------
\usepackage[T1]{fontenc}
\usepackage{textcomp} % Euro-Zeichen etc.

% Schrift ----------------------------------------------------------------------
\usepackage{lmodern} % bessere Fonts
\usepackage{relsize} % Schriftgröße relativ festlegen

% Tabellen ---------------------------------------------------------------------
\PassOptionsToPackage{table}{xcolor}
\usepackage{tabularx}
% für lange Tabellen
\usepackage{longtable}
\usepackage{array}
\usepackage{ragged2e}
\usepackage{lscape}
\newcolumntype{w}[1]{>{\raggedleft\hspace{0pt}}p{#1}} % Spaltendefinition rechtsbündig mit definierter Breite

% Grafiken ---------------------------------------------------------------------
\usepackage[dvips,final]{graphicx} % Einbinden von JPG-Grafiken ermöglichen
\usepackage{graphics} % keepaspectratio
\usepackage{floatflt} % zum Umfließen von Bildern
\graphicspath{{Bilder/}} % hier liegen die Bilder des Dokuments

% Sonstiges --------------------------------------------------------------------
\usepackage[titles]{tocloft} % Inhaltsverzeichnis DIN 5008 gerecht einrücken
\usepackage{amsmath,amsfonts} % Befehle aus AMSTeX für mathematische Symbole
\usepackage{enumitem} % anpassbare Enumerates/Itemizes
\usepackage{xspace} % sorgt dafür, dass Leerzeichen hinter parameterlosen Makros nicht als Makroendezeichen interpretiert werden

\usepackage{makeidx} % für Index-Ausgabe mit \printindex
\usepackage[printonlyused]{acronym} % es werden nur benutzte Definitionen aufgelistet

% Einfache Definition der Zeilenabstände und Seitenränder etc.
\usepackage{setspace}
\usepackage{geometry}

% Symbolverzeichnis
\usepackage[intoc]{nomencl}
\let\abbrev\nomenclature
\renewcommand{\nomname}{Abkürzungsverzeichnis}
\setlength{\nomlabelwidth}{.25\hsize}
\renewcommand{\nomlabel}[1]{#1 \dotfill}
\setlength{\nomitemsep}{-\parsep}

\usepackage{varioref} % Elegantere Verweise. „auf der nächsten Seite“
\usepackage{url} % URL verlinken, lange URLs umbrechen etc.

\usepackage{chngcntr} % fortlaufendes Durchnummerieren der Fußnoten
% \usepackage[perpage]{footmisc} % Alternative: Nummerierung der Fußnoten auf jeder Seite neu

\usepackage{ifthen} % bei der Definition eigener Befehle benötigt
\usepackage{todonotes} % definiert u.a. die Befehle \todo und \listoftodos
\usepackage[square]{natbib} % wichtig für korrekte Zitierweise

% PDF-Optionen -----------------------------------------------------------------
\usepackage{pdfpages}
\pdfminorversion=5 % erlaubt das Einfügen von pdf-Dateien bis Version 1.7, ohne eine Fehlermeldung zu werfen (keine Garantie für fehlerfreies Einbetten!)
\usepackage[
    bookmarks,
    bookmarksnumbered,
    bookmarksopen=true,
    bookmarksopenlevel=1,
    colorlinks=true,
% diese Farbdefinitionen zeichnen Links im PDF farblich aus
    linkcolor=AOBlau, % einfache interne Verknüpfungen
    anchorcolor=AOBlau,% Ankertext
    citecolor=AOBlau, % Verweise auf Literaturverzeichniseinträge im Text
    filecolor=AOBlau, % Verknüpfungen, die lokale Dateien öffnen
    menucolor=AOBlau, % Acrobat-Menüpunkte
    urlcolor=AOBlau,
% diese Farbdefinitionen sollten für den Druck verwendet werden (alles schwarz)
    %linkcolor=black, % einfache interne Verknüpfungen
    %anchorcolor=black, % Ankertext
    %citecolor=black, % Verweise auf Literaturverzeichniseinträge im Text
    %filecolor=black, % Verknüpfungen, die lokale Dateien öffnen
    %menucolor=black, % Acrobat-Menüpunkte
    %urlcolor=black,
%
    %backref, % Quellen werden zurück auf ihre Zitate verlinkt
    pdftex,
    plainpages=false, % zur korrekten Erstellung der Bookmarks
    pdfpagelabels=true, % zur korrekten Erstellung der Bookmarks
    hypertexnames=false, % zur korrekten Erstellung der Bookmarks
    linktocpage % Seitenzahlen anstatt Text im Inhaltsverzeichnis verlinken
]{hyperref}
% Befehle, die Umlaute ausgeben, führen zu Fehlern, wenn sie hyperref als Optionen übergeben werden
\hypersetup{
    pdftitle={\titel -- \untertitel},
    pdfauthor={\autorName},
    pdfcreator={\autorName},
    pdfsubject={\titel -- \untertitel},
    pdfkeywords={\titel -- \untertitel},
}


% zum Einbinden von Programmcode -----------------------------------------------
\usepackage{listings}
\usepackage{xcolor}
\definecolor{hellgelb}{rgb}{1,1,0.9}
\definecolor{colKeys}{rgb}{0,0,1}
\definecolor{colIdentifier}{rgb}{0,0,0}
\definecolor{colComments}{rgb}{0,0.5,0}
\definecolor{colString}{rgb}{1,0,0}
\lstset{
    float=hbp,
	basicstyle=\footnotesize,
    identifierstyle=\color{colIdentifier},
    keywordstyle=\color{colKeys},
    stringstyle=\color{colString},
    commentstyle=\color{colComments},
    backgroundcolor=\color{hellgelb},
    columns=flexible,
    tabsize=2,
    frame=single,
    extendedchars=true,
    showspaces=false,
    showstringspaces=false,
    numbers=left,
    numberstyle=\tiny,
    breaklines=true,
    breakautoindent=true,
	captionpos=b,
}
\lstdefinelanguage{cs}{
	sensitive=false,
	morecomment=[l]{//},
	morecomment=[s]{/*}{*/},
	morestring=[b]",
	morekeywords={
		abstract,event,new,struct,as,explicit,null,switch
		base,extern,object,this,bool,false,operator,throw,
		break,finally,out,true,byte,fixed,override,try,
		case,float,params,typeof,catch,for,private,uint,
		char,foreach,protected,ulong,checked,goto,public,unchecked,
		class,if,readonly,unsafe,const,implicit,ref,ushort,
		continue,in,return,using,decimal,int,sbyte,virtual,
		default,interface,sealed,volatile,delegate,internal,short,void,
		do,is,sizeof,while,double,lock,stackalloc,
		else,long,static,enum,namespace,string},
}
\lstdefinelanguage{natural}{
	sensitive=false,
	morecomment=[l]{/*},
	morestring=[b]",
	morestring=[b]',
	alsodigit={-,*},
	morekeywords={
		DEFINE,DATA,LOCAL,END-DEFINE,WRITE,CALLNAT,PARAMETER,USING,
		IF,NOT,END-IF,ON,*ERROR-NR,ERROR,END-ERROR,ESCAPE,ROUTINE,
		PERFORM,SUBROUTINE,END-SUBROUTINE,CONST,END-FOR,END,FOR,RESIZE,
		ARRAY,TO,BY,VALUE,RESET,COMPRESS,INTO,EQ},
}
\lstdefinelanguage{php}{
	sensitive=false,
	morecomment=[l]{/*},
	morestring=[b]",
	morestring=[b]',
	alsodigit={-,*},
	morekeywords={
		abstract,and,array,as,break,case,catch,cfunction,class,clone,const,
		continue,declare,default,do,else,elseif,enddeclare,endfor,endforeach,
		endif,endswitch,endwhile,extends,final,for,foreach,function,global,
		goto,if,implements,interface,instanceof,namespace,new,old_function,or,
		private,protected,public,static,switch,throw,try,use,var,while,xor
		die,echo,empty,exit,eval,include,include_once,isset,list,require,
		require_once,return,print,unset},
}
 % verwendete Packages
% !TEX root = ../Projektdokumentation.tex

% Seitenränder -----------------------------------------------------------------
\setlength{\topskip}{\ht\strutbox} % behebt Warnung von geometry
\geometry{a4paper,left=20mm,right=20mm,top=25mm,bottom=35mm}

\usepackage[
	automark, % Kapitelangaben in Kopfzeile automatisch erstellen
	headsepline, % Trennlinie unter Kopfzeile
	ilines % Trennlinie linksbündig ausrichten
]{scrpage2}

% Kopf- und Fußzeilen ----------------------------------------------------------
\pagestyle{scrheadings}
% chapterpagestyle gibt es nicht in scrartcl
%\renewcommand{\chapterpagestyle}{scrheadings}
\clearscrheadfoot

% Kopfzeile
\renewcommand{\headfont}{\normalfont} % Schriftform der Kopfzeile
\ihead{\large{\textsc{\titel}}\\ \small{\untertitel} \\[2ex] \textit{\headmark}}
\chead{}
\ohead{\includegraphics[scale=0.3]{\betriebLogo}}
\setlength{\headheight}{15mm} % Höhe der Kopfzeile
%\setheadwidth[0pt]{textwithmarginpar} % Kopfzeile über den Text hinaus verbreitern (falls Logo den Text überdeckt)

% Fußzeile
\ifoot{\autorName}
\cfoot{}
\ofoot{\pagemark}

% Überschriften nach DIN 5008 in einer Fluchtlinie
% ------------------------------------------------------------------------------

% Abstand zwischen Nummerierung und Überschrift definieren
% > Schön wäre hier die dynamische Berechnung des Abstandes in Abhängigkeit
% > der Verschachtelungstiefe des Inhaltsverzeichnisses
\newcommand{\headingSpace}{1.5cm}

% Abschnittsüberschriften im selben Stil wie beim Inhaltsverzeichnis einrücken
\renewcommand*{\othersectionlevelsformat}[3]{
  \makebox[\headingSpace][l]{#3\autodot}
}

% Für die Einrückung wird das Paket tocloft benötigt
%\cftsetindents{chapter}{0.0cm}{\headingSpace}
\cftsetindents{section}{0.0cm}{\headingSpace}
\cftsetindents{subsection}{0.0cm}{\headingSpace}
\cftsetindents{subsubsection}{0.0cm}{\headingSpace}
\cftsetindents{figure}{0.0cm}{\headingSpace}
\cftsetindents{table}{0.0cm}{\headingSpace}


% Allgemeines
% ------------------------------------------------------------------------------

\onehalfspacing % Zeilenabstand 1,5 Zeilen
\frenchspacing % erzeugt ein wenig mehr Platz hinter einem Punkt

% Schusterjungen und Hurenkinder vermeiden
\clubpenalty = 10000
\widowpenalty = 10000
\displaywidowpenalty = 10000

% Quellcode-Ausgabe formatieren
\lstset{numbers=left, numberstyle=\tiny, numbersep=5pt, breaklines=true}
\lstset{emph={square}, emphstyle=\color{red}, emph={[2]root,base}, emphstyle={[2]\color{blue}}}

\counterwithout{footnote}{section} % Fußnoten fortlaufend durchnummerieren
\setcounter{tocdepth}{\subsubsectionlevel} % im Inhaltsverzeichnis werden die Kapitel bis zum Level der subsubsection übernommen
\setcounter{secnumdepth}{\subsubsectionlevel} % Kapitel bis zum Level der subsubsection werden nummeriert

% Aufzählungen anpassen
\renewcommand{\labelenumi}{\arabic{enumi}.}
\renewcommand{\labelenumii}{\arabic{enumi}.\arabic{enumii}.}
\renewcommand{\labelenumiii}{\arabic{enumi}.\arabic{enumii}.\arabic{enumiii}}

% Tabellenfärbung:
\definecolor{heading}{rgb}{0.64,0.78,0.86}
\definecolor{odd}{rgb}{0.9,0.9,0.9}
 % Definitionen zum Aussehen der Seiten
% !TEX root = ../Projektdokumentation.tex

% Abkürzungen, ggfs. mit korrektem Leerraum
\newcommand{\bs}{$\backslash$\xspace}
\newcommand{\bspw}{bspw.\xspace}
\newcommand{\bzw}{bzw.\xspace}
\newcommand{\ca}{ca.\xspace}
\newcommand{\dahe}{\mbox{d.\,h.}\xspace}
\newcommand{\etc}{etc.\xspace}
\newcommand{\eur}[1]{\mbox{#1\,\texteuro}\xspace}
\newcommand{\evtl}{evtl.\xspace}
\newcommand{\ggfs}{ggfs.\xspace}
\newcommand{\Ggfs}{Ggfs.\xspace}
\newcommand{\gqq}[1]{\glqq{}#1\grqq{}}
\newcommand{\inkl}{inkl.\xspace}
\newcommand{\insb}{insb.\xspace}
\newcommand{\ua}{\mbox{u.\,a.}\xspace}
\newcommand{\usw}{usw.\xspace}
\newcommand{\Vgl}{Vgl.\xspace}
\newcommand{\zB}{\mbox{z.\,B.}\xspace}

% Befehle für häufig anfallende Aufgaben
\newcommand{\Abbildung}[1]{\autoref{fig:#1}}
\newcommand{\Anhang}[1]{\appendixname{}~\ref{#1}: \nameref{#1} \vpageref{#1}}
\newcommand{\includegraphicsKeepAspectRatio}[2]{\includegraphics[width=#2\textwidth,height=#2\textheight,keepaspectratio]{#1}}
\newcommand{\Zitat}[2][\empty]{\ifthenelse{\equal{#1}{\empty}}{\citep{#2}}{\citep[#1]{#2}}}
\newcommand{\Autor}[1]{\textsc{#1}} % zum Ausgeben von Autoren
\newcommand{\itemd}[2]{\item{\textbf{#1}}{#2}} % erzeugt ein Listenelement mit fetter Überschrift

% fügt Tabellen aus einer TEX-Datei ein
\newcommand{\tabelle}[3] % Parameter: caption, label, file
{\begin{table}[htbp]
\centering
\singlespacing
\input{Tabellen/#3}
\caption{#1}
\label{#2}
\end{table}}

\newcommand{\tabelleAnhang}[1] % Parameter: file
{\begin{center}
\singlespacing
\input{Tabellen/#1}
\end{center}}

% einfaches Wechseln der Schrift, z.B.: \changefont{cmss}{sbc}{n}
\newcommand{\changefont}[3]{\fontfamily{#1} \fontseries{#2} \fontshape{#3} \selectfont}

% Verwendung analog zu \includegraphics
\newlength{\myx} % Variable zum Speichern der Bildbreite
\newlength{\myy} % Variable zum Speichern der Bildhöhe
\newcommand\includegraphicstotab[2][\relax]{%
% Abspeichern der Bildabmessungen
\settowidth{\myx}{\includegraphics[{#1}]{#2}}%
\settoheight{\myy}{\includegraphics[{#1}]{#2}}%
% das eigentliche Einfügen
\parbox[c][1.1\myy][c]{\myx}{%
\includegraphics[{#1}]{#2}}%
}

\definecolor{AOBlau}{rgb}{0, 0.28, 0.56}

% verschiedene Befehle um Wörter semantisch auszuzeichnen ----------------------
\newcommand{\Index}[2][\empty]{\ifthenelse{\equal{#1}{\empty}}{\index{#2}#2}{\index{#1}#2}}
\newcommand{\Fachbegriff}[2][\empty]{\ifthenelse{\equal{#1}{\empty}}{\textit{\Index{#2}}}{\textit{\Index[#1]{#2}}}}
\newcommand{\NeuerBegriff}[2][\empty]{\ifthenelse{\equal{#1}{\empty}}{\textbf{\Index{#2}}}{\textbf{\Index[#1]{#2}}}}

\newcommand{\Ausgabe}[1]{\texttt{#1}}
\newcommand{\Eingabe}[1]{\texttt{#1}}
\newcommand{\Code}[1]{\texttt{#1}}
\newcommand{\Datei}[1]{\texttt{#1}}

\newcommand{\Assembly}[1]{\textsf{#1}}
\newcommand{\Klasse}[1]{\textsf{#1}}
\newcommand{\Methode}[1]{\textsf{#1}}
\newcommand{\Attribut}[1]{\textsf{#1}}

\newcommand{\Datentyp}[1]{\textsf{#1}}
\newcommand{\XMLElement}[1]{\textsf{#1}}
\newcommand{\Webservice}[1]{\textsf{#1}}

\newcommand{\Refactoring}[1]{\Fachbegriff{#1}}
\newcommand{\CodeSmell}[1]{\Fachbegriff{#1}}
\newcommand{\Metrik}[1]{\Fachbegriff{#1}}
\newcommand{\DesignPattern}[1]{\Fachbegriff{#1}}
 % eigene allgemeine Befehle, die z.B. die Arbeit mit LaTeX erleichtern
% !TEX root = ../Projektdokumentation.tex

% Abkürzungen, ggfs. mit korrektem Leerraum
\newcommand{\bs}{$\backslash$\xspace}
\newcommand{\bspw}{bspw.\xspace}
\newcommand{\bzw}{bzw.\xspace}
\newcommand{\ca}{ca.\xspace}
\newcommand{\dahe}{\mbox{d.\,h.}\xspace}
\newcommand{\etc}{etc.\xspace}
\newcommand{\eur}[1]{\mbox{#1\,\texteuro}\xspace}
\newcommand{\evtl}{evtl.\xspace}
\newcommand{\ggfs}{ggfs.\xspace}
\newcommand{\Ggfs}{Ggfs.\xspace}
\newcommand{\gqq}[1]{\glqq{}#1\grqq{}}
\newcommand{\inkl}{inkl.\xspace}
\newcommand{\insb}{insb.\xspace}
\newcommand{\ua}{\mbox{u.\,a.}\xspace}
\newcommand{\usw}{usw.\xspace}
\newcommand{\Vgl}{Vgl.\xspace}
\newcommand{\zB}{\mbox{z.\,B.}\xspace}

% Befehle für häufig anfallende Aufgaben
\newcommand{\Abbildung}[1]{\autoref{fig:#1}}
\newcommand{\Anhang}[1]{\appendixname{}~\ref{#1}: \nameref{#1} \vpageref{#1}}
\newcommand{\includegraphicsKeepAspectRatio}[2]{\includegraphics[width=#2\textwidth,height=#2\textheight,keepaspectratio]{#1}}
\newcommand{\Zitat}[2][\empty]{\ifthenelse{\equal{#1}{\empty}}{\citep{#2}}{\citep[#1]{#2}}}
\newcommand{\Autor}[1]{\textsc{#1}} % zum Ausgeben von Autoren
\newcommand{\itemd}[2]{\item{\textbf{#1}}{#2}} % erzeugt ein Listenelement mit fetter Überschrift

% fügt Tabellen aus einer TEX-Datei ein
\newcommand{\tabelle}[3] % Parameter: caption, label, file
{\begin{table}[htbp]
\centering
\singlespacing
\input{Tabellen/#3}
\caption{#1}
\label{#2}
\end{table}}

\newcommand{\tabelleAnhang}[1] % Parameter: file
{\begin{center}
\singlespacing
\input{Tabellen/#1}
\end{center}}

% einfaches Wechseln der Schrift, z.B.: \changefont{cmss}{sbc}{n}
\newcommand{\changefont}[3]{\fontfamily{#1} \fontseries{#2} \fontshape{#3} \selectfont}

% Verwendung analog zu \includegraphics
\newlength{\myx} % Variable zum Speichern der Bildbreite
\newlength{\myy} % Variable zum Speichern der Bildhöhe
\newcommand\includegraphicstotab[2][\relax]{%
% Abspeichern der Bildabmessungen
\settowidth{\myx}{\includegraphics[{#1}]{#2}}%
\settoheight{\myy}{\includegraphics[{#1}]{#2}}%
% das eigentliche Einfügen
\parbox[c][1.1\myy][c]{\myx}{%
\includegraphics[{#1}]{#2}}%
}

\definecolor{AOBlau}{rgb}{0, 0.28, 0.56}

% verschiedene Befehle um Wörter semantisch auszuzeichnen ----------------------
\newcommand{\Index}[2][\empty]{\ifthenelse{\equal{#1}{\empty}}{\index{#2}#2}{\index{#1}#2}}
\newcommand{\Fachbegriff}[2][\empty]{\ifthenelse{\equal{#1}{\empty}}{\textit{\Index{#2}}}{\textit{\Index[#1]{#2}}}}
\newcommand{\NeuerBegriff}[2][\empty]{\ifthenelse{\equal{#1}{\empty}}{\textbf{\Index{#2}}}{\textbf{\Index[#1]{#2}}}}

\newcommand{\Ausgabe}[1]{\texttt{#1}}
\newcommand{\Eingabe}[1]{\texttt{#1}}
\newcommand{\Code}[1]{\texttt{#1}}
\newcommand{\Datei}[1]{\texttt{#1}}

\newcommand{\Assembly}[1]{\textsf{#1}}
\newcommand{\Klasse}[1]{\textsf{#1}}
\newcommand{\Methode}[1]{\textsf{#1}}
\newcommand{\Attribut}[1]{\textsf{#1}}

\newcommand{\Datentyp}[1]{\textsf{#1}}
\newcommand{\XMLElement}[1]{\textsf{#1}}
\newcommand{\Webservice}[1]{\textsf{#1}}

\newcommand{\Refactoring}[1]{\Fachbegriff{#1}}
\newcommand{\CodeSmell}[1]{\Fachbegriff{#1}}
\newcommand{\Metrik}[1]{\Fachbegriff{#1}}
\newcommand{\DesignPattern}[1]{\Fachbegriff{#1}}
 % eigene projektspezifische Befehle, z.B. Abkürzungen usw.

\begin{document}

% kann nach dem Lesen entfernt werden ---------------------------------------
%\pagestyle{plain}
%% !TEX root = Projektdokumentation.tex
\section*{Über diese Vorlage}

Diese \LaTeX-Vorlage wurde von Stefan Macke\footnote{Blog des Autors:
\url{http://fachinformatiker-anwendungsentwicklung.net}, Twitter:
\Eingabe{@StefanMacke}} als Grundlage für die Projektdokumentationen der Auszubildenden zum Fachinformatiker mit Fachrichtung
Anwendungsentwicklung bei der \AO entwickelt. Nichtsdestotrotz dürfte sie ebenso für die anderen IT-Berufe\footnote{\zB IT-Kaufleute, Fachinformatiker
mit Fachrichtung Systemintegration \usw} geeignet sein, da diese anhand der gleichen Verordnung bewertet werden.

Diese Vorlage enthält bereits eine Vorstrukturierung der möglichen Inhalte einer tatsächlichen Projektdokumentation, die auf Basis der
Erfahrungen im Rahmen der Prüfertätigkeit des Autors erstellt und unter Zuhilfenahme von \citet{Rohrer2011} abgerundet wurden.

Sämtliche verwendeten Abbildungen, Tabellen und Listings stammen von \citet{Grashorn2010}.

Download-Link für diese Vorlage: \url{http://fiae.link/LaTeXVorlageFIAE}

Auch verfügbar auf GitHub: \url{https://github.com/StefanMacke/latex-vorlage-fiae}

\subsection*{Lizenz}

\begin{center}
\includegraphicsKeepAspectRatio{CC-Logo.pdf}{0.3}
\end{center}
Dieses Werk steht unter einer Creative Commons Namensnennung - Weitergabe unter gleichen Bedingungen 4.0 International Lizenz.
\footnote{\url{http://creativecommons.org/licenses/by-sa/4.0/}}

\begin{center}
\includegraphicsKeepAspectRatio{CC-Attribution.pdf}{0.07}
\includegraphicsKeepAspectRatio{CC-ShareAlike.pdf}{0.07}
\end{center}

\begin{description}
	\item[Namensnennung] Sie müssen den Namen des Autors/Rechteinhabers in der von ihm festgelegten Weise nennen.
	\footnote{Die Namensnennung im \LaTeX-Quelltext mit Link auf \url{http://fiae.link/LaTeXVorlageFIAE} reicht hierfür aus.}
	\item[Weitergabe unter gleichen Bedingungen] Wenn Sie das lizenzierte Werk \bzw den lizenzierten Inhalt bearbeiten
	oder in anderer Weise erkennbar als Grundlage für eigenes Schaffen verwenden, dürfen Sie die daraufhin neu entstandenen
	Werke \bzw Inhalte nur unter Verwendung von Lizenzbedingungen weitergeben, die mit denen dieses Lizenzvertrages identisch oder vergleichbar sind.
\end{description}

\subsection*{Inhalt der Projektdokumentation}

Grundsätzlich definiert die \citet[S.~1746]{Bundesgesetzblatt48}\footnote{Dieses
Dokument sowie alle weiteren hier genannten können unter
\url{http://fiae.link/LaTeXVorlageFIAEQuellen} heruntergeladen werden.} das Ziel der Projektdokumentation wie folgt:
\begin{quote}
"`Durch die Projektarbeit und deren Dokumentation soll der Prüfling belegen, daß er Arbeitsabläufe und Teilaufgaben zielorientiert unter
Beachtung wirtschaftlicher, technischer, organisatorischer und zeitlicher Vorgaben selbständig planen und kundengerecht umsetzen sowie
Dokumentationen kundengerecht anfertigen, zusammenstellen und modifizieren kann."'
\end{quote}

Und das \citet[S.~36]{BMBF2000} ergänzt:
\begin{quote}
"`Die Ausführung der Projektarbeit wird mit praxisbezogenen Unterlagen dokumentiert.
Der Prüfungsausschuss bewertet die Projektarbeit anhand der Dokumentation. Dabei
wird nicht das Ergebnis -- \zB ein lauffähiges Programm -- herangezogen, sondern
der Arbeitsprozess. Die Dokumentation ist keine wissenschaftliche Abhandlung,
sondern eine handlungsorientierte Darstellung des Projektablaufs mit
praxisbezogenen, d.h. betriebüblichen Unterlagen. Sie soll einen Umfang von
maximal 10 bis 15 DIN A 4-Seiten nicht überschreiten. Soweit erforderlich können in
einem Anhang \zB den Zusammenhang erläuternde Darstellungen beigefügt werden."'
\end{quote}

Außerdem werden dort die grundlegenden Inhalte der Projektdokumentation aufgelistet:
\begin{itemize}
	\item Name und Ausbildungsberuf des Prüfungsteilnehmers
	\item Angabe des Ausbildungsbetriebes
	\item Thema der Projektarbeit
	\item Falls erforderlich, Beschreibung/Konkretisierung des Auftrages
	\item Umfassende Beschreibung der Prozessschritte und der erzielten Ergebnisse
	\item Gegebenenfalls Veränderungen zum Projektantrag mit Begründung
	\item Wenn für das Projekt erforderlich, ein Anhang mit praxisbezogenen Unterlagen und Dokumenten. Dieser Anhang sollte nicht
	aufgebläht werden. Die angehängten Dokumente und Unterlagen sind auf das absolute Minimum zu beschränken.
\end{itemize}

In den folgenden Kapiteln werden diese geforderten Inhalte und sinnvolle Ergänzungen nun meist stichwortartig und \ggfs mit
Beispielen beschrieben. Nicht alle Kapitel müssen in jeder Dokumentation vorhanden sein. Handelt es sich \bspw um ein in sich
geschlossenes Projekt, kann das Kapitel~\ref{sec:Projektabgrenzung}: \nameref{sec:Projektabgrenzung} entfallen; arbeitet die
Anwendung nur mit \acs{XML}-Dateien, kann und muss keine Datenbank beschrieben werden \usw


\subsection*{Formale Vorgaben}

Die formalen Vorgaben zum Umfang und zur Gestaltung der Projektdokumentation können je nach IHK recht unterschiedlich sein.
Normalerweise sollte die zuständige IHK einen Leitfaden bereitstellen, in dem alle Formalien nachgelesen werden können,
wie \zB bei der \citet{MerkblattIHK}.

Als Richtwert verwende ich 15 Seiten für den reinen Inhalt. Also in dieser Vorlage alle Seiten, die arabisch nummeriert
sind (ohne das Literaturverzeichnis und die eidesstattliche Erklärung).
Große Abbildungen, Quelltexte, Tabellen \usw gehören in den Anhang, der 25 Seiten nicht überschreiten sollte.

Typographische Konventionen, Seitenränder \usw können in der Datei \Datei{Seitenstil.tex} beliebig angepasst werden.


\subsection*{Bewertungskriterien}
Die Bewertungskriterien für die Benotung der Projektdokumentation sind recht einheitlich und können leicht in Erfahrung
gebracht werden, \zB bei der \citet{BewertungsmatrikIHK}.
Grundsätzlich sollte die Projektdokumentation nach der Fertigstellung noch einmal im Hinblick auf diese Kriterien durchgeschaut werden.

%\cleardoublepage
%%\pagestyle{scrheadings}
% ---------------------------------------------------------------------------

% Deckblatt -----------------------------------------------------------------
\phantomsection
\thispagestyle{empty}
\pdfbookmark[1]{Eidesstattliche Erklärung}{ihkdeckblatt}
\includegraphicsKeepAspectRatio{DeckblattIHK}{1}
\cleardoublepage

\phantomsection
\thispagestyle{plain}
\pdfbookmark[1]{Deckblatt}{deckblatt}
% !TEX root = Projektdokumentation.tex
\begin{titlepage}

\begin{center}
\includegraphics[scale=0.25]{LogoIHK.pdf}\\[1ex]
\Large{Abschlussprüfung \pruefungstermin}\\[3ex]

\Large{\ausbildungsberuf}\\
\LARGE{\betreff}\\[4ex]

\huge{\textbf{\titel}}\\[1.5ex]
\Large{\textbf{\untertitel}}\\[4ex]

\normalsize
Abgabetermin: \abgabeOrt, den \abgabeTermin\\[3em]
\textbf{Prüfungsbewerber:}\\
\autorName\\
\autorAnschrift\\
\autorOrt\\[5ex]

\includegraphics[scale=1]{\betriebLogoMitTitel}\\[2ex]
\textbf{Ausbildungsbetrieb:}\\
\betriebName\\
\betriebAnschrift\\
\betriebOrt\\[5em]
\end{center}

\small
\noindent
Dieses Werk einschließlich seiner Teile ist \textbf{urheberrechtlich geschützt}.
Jede Verwertung außerhalb der engen Grenzen des Urheberrechtgesetzes ist ohne
Zustimmung des Autors unzulässig und strafbar. Das gilt insbesondere für
Vervielfältigungen, Übersetzungen, Mikroverfilmungen sowie die Einspeicherung
und Verarbeitung in elektronischen Systemen.

\end{titlepage}
\cleardoublepage

% Preface --------------------------------------------------------------------
\phantomsection
\pagenumbering{Roman}
\pdfbookmark[1]{Inhaltsverzeichnis}{inhalt}
\tableofcontents
\cleardoublepage

\phantomsection
\listoffigures
\cleardoublepage

\phantomsection
\listoftables
\cleardoublepage

%\phantomsection
%\lstlistoflistings
%\cleardoublepage

\newcommand{\abkvz}{Abkürzungsverzeichnis}
\renewcommand{\nomname}{\abkvz}
\section*{\abkvz}
\markboth{\abkvz}{\abkvz}
\addcontentsline{toc}{section}{\abkvz}
% !TEX root = Projektdokumentation.tex

% Es werden nur die Abkürzungen aufgelistet, die mit \ac definiert und auch benutzt wurden. 
%
% \acro{VERSIS}{Versicherungsinformationssystem\acroextra{ (Bestandsführungssystem)}}
% Ergibt in der Liste: VERSIS Versicherungsinformationssystem (Bestandsführungssystem)
% Im Text aber: \ac{VERSIS} -> Versicherungsinformationssystem (VERSIS)

% Hinweis: allgemein bekannte Abkürzungen wie z.B. bzw. u.a. müssen nicht ins Abkürzungsverzeichnis aufgenommen werden
% Hinweis: allgemein bekannte IT-Begriffe wie Datenbank oder Programmiersprache müssen nicht erläutert werden,
%          aber ggfs. Fachbegriffe aus der Domäne des Prüflings (z.B. Versicherung)

% Die Option (in den eckigen Klammern) enthält das längste Label oder
% einen Platzhalter der die Breite der linken Spalte bestimmt.
\begin{acronym}[WWWWWW]
	\acro{API}{Application Programming Interface}
	\acro{HTTP}{Hypertext Transport Protocol}
	\acro{SDK}{Software Development Kit}
	\acro{ERM}{En\-ti\-ty-Re\-la\-tion\-ship-Mo\-dell}
	\acro{UML}{Unified Modeling Language}
	\acro{GUI}{Grafische Benutzeroberfläche}
	\acro{SVN}{Subversion}
	
	\acro{DRV}{\betriebName}
	\acro{CBP-AD}{Computergestütze Betriebsprüfung - Abschluss \& Dokumente}
	\acro{CBP-NB}{Computergestütze Betriebsprüfung - Nachberechnung}
	
	\acro{IDE}{Integrated Development Environment}
	\acro{MVC}[MVC]{Model View Controller}
	\acro{ORM}{Object-Relational Mapping}
	\acro{PHP}{Hypertext Preprocessor}
	\acro{SQL}{Structured Query Language}
	\acro{EPK}{Ereignisgesteuerte Prozesskette}
	
	\acro{XML}{Extensible Markup Language}
	\acro{HTML}{Hypertext Markup Language}\acused{HTML}
	\acro{CSS}{Cascading Style Sheets}
	\acro{CSV}{Comma Separated Value}
	
	\acro{AO}{\textsc{Alte Oldenburger} Krankenversicherung AG}
	\acro{NatInfo}[\textsc{NatInfo}]{Natural Information System}
	\acro{Natural}[\textsc{Natural}]{Programmiersprache der Software AG}\acused{Natural}
\end{acronym}

\clearpage

% Inhalt ---------------------------------------------------------------------
\pagenumbering{arabic}
% !TEX root = Projektdokumentation.tex
% !TEX root = ../Projektdokumentation.tex

\section{Einleitung}
\label{sec:Einleitung}


\subsection{Projektumfeld} 
\label{sec:Projektumfeld}

Die \ac{DRV Bund} ist ein bundesweit tätiger Träger der gesetzlichen Rentenversicherung in der Bundesrepublik Deutschland
mit \ca 17000 Mitarbeitern.\\
Zum Aufgabenfeld gehören:
\begin{itemize}
	\item Bearbeitung von Rentenanträgen und Auszahlung von Renten
	\item Überprüfung von Sozialabgaben auf Richtigkeit
	\item Beratung zu gesetzlichen Pflichten und Rechten
	\item Bearbeitung von Rehabilitationsanträgen und Beaufsichtigung der Rehabilitationseinrichtungen
\end{itemize}

Für die Überprüfung von Sozialabgaben entwickelt die IT-Abteilung der \DRV verschiedene Anwendung, um diesen Prozess zu vereinfachen.

Die Betriebsprüfung

Die \ac{CBP-AD} ist eine Desktopanwendung mit der jeweilig benötigte Dokumente, Anlagen und Schreiben(nachfolgend nur noch Dokumente genannt) für Betriebsprüfungen erzeugt werden können. Diese Dokumente basieren auf Datenquellen der Desktopanwendung \ac{CBP-NB}.


\subsection{Projektziel} 
\label{sec:Projektziel}

Die Betriebsprüfer erstellen und bearbeiten Dokumente und Anlagen, die mit Daten aus Berechnungen der Anwendung \acs{CBP-NB} befüllt werden. Wenn in dieser Daten verändert werden, die schon in erzeugten Dokumenten vorkommen, sind diese Dokumente in einem „asynchronen“ Zustand und müssen vor Weiterverwendung aktualisiert werden.

Im Hauptmenü der \acs{CBP-AD} soll ein neuer Menüeintrag bereitgestellt werden, dessen Kommando einen Aktualisierungsprozess anstößt, der alle Dokumente auf Asynchronität prüft und anschließend veraltete aktualisiert.

Für die verschiedenen Dokumenttypen gibt es zurzeit auch noch unterschiedliche Strukturen, wie die jeweiligen Dokumente neu erzeugt werden. Für die Funktionalität „Globales Aktualisieren“ sollen nun alle Dokumenttypen auf eine Struktur umgebaut werden.

Erzeugungsstrukturen für einige Dokumenttypen berechnen ihren Fortschritt eigenständig und geben diesen in einer eigenen Oberfläche aus. Für diese soll eine Möglichkeit der Unterdrückung dieser Fortschrittsausgabe implementiert werden, damit der Aktualisierungsprozess „Globales Aktualisieren“ dies einheitlich für alle Dokumenttypen ausgeben kann.


\subsection{Projektbegründung} 
\label{sec:Projektbegruendung}

Durch diese Erweiterung wird eine Vereinheitlichung der Dokumentenaktualisierung erreicht, die zugleich eine erhebliche Vereinfachung für den Anwender mit sich bringt.


\subsection{Projektschnittstellen} 
\label{sec:Projektschnittstellen}

Daten aus den Berechnungen der Sozialabgaben von Betrieben und ihren Mitarbeitern werden über eine \acs{HTTP}-\acs{API} der \acs{CBP-NB} angefordert.

Die Benutzer der Anwendung sind die Betriebsprüfer der \DRV.


\subsection{Projektabgrenzung} 
\label{sec:Projektabgrenzung}

Dieses Projekt zur Erweiterung der \acs{CBP-AD} ist unabhängig von der Entwicklung der \ac{CBP-NB}, da nur bereits festgelegte Schnittstellen zum Datenaustausch benutzt werden und keine Änderung dieser notwendig sind.
% !TEX root = ../Projektdokumentation.tex


\section{Projektplanung} 
\label{sec:Projektplanung}


\subsection{Projektphasen}
\label{sec:Projektphasen}

Die (Vorbereitungs-)Projektphase begann am 17.10.2016 und endete am 16.12.2016.
Die tägliche Arbeitszeit betrug 7 Stunden 48 Minuten mit 30 Minuten Mittagspause.
Die Projektarbeit fand nicht durchgängig statt, da Betriebsinterne Ereignis berücksichtigt werden müssen.

\tabelle{Zeitplanung}{tab:Zeitplanung}{Zeitplanung}

Eine detailliertere Zeitplanung findet sich im \Anhang{app:Zeitplanung}.

\subsection{Abweichungen vom Projektantrag}
\label{sec:AbweichungenProjektantrag}

Entwurfsphase mit 5 Stunden ist dazu gekommen, da eine Planung her muss.

Implementierungsphase ist daher um gleiche Zeitspanne verringert.  % AAbweichungen vervollständigen

\subsection{Ressourcenplanung}
\label{sec:Ressourcenplanung}

\begin{itemize}
\renewcommand{\labelitemi}{}
\renewcommand{\labelitemii}{\textbullet}
\renewcommand{\labelitemiii}{\normalfont\bfseries\textendash}
	\itemd{Hardware}{
		\begin{itemize}
			\item Büroarbeitsplatz mit Tisch, Stuhl, Stromanschlüsse
			\item Arbeitsmaschine 1 mit Windows7
			\item Arbeitsmaschine 2 mit Entwicklungsnetz-Zugang + Kartenleser
		\end{itemize}
	}
	\itemd{Software}{
		\begin{itemize}
			\item Visual Studio Professional 2013 + .NET-Framework (mindestens v2.0)
			\item \ac{SDK}s für Visual Studio
			\begin{itemize}
				\itemd{CrytalReports for VisualStudio}{Framework zum Erzeugung von PDFs aus Dokumenten}
				\itemd{GhostScript}{Framework zum Drucken von Dokumenten}
				\itemd{TX-Spell-dotNet-0300}{Framework zur Rechtschreibprüfung von Dokumenten}		\itemd{TX-Text-Control-dotNet-1900}{Framework zur Darstellung von Dokumenten}
			\end{itemize}
			\item MiKTeX - Distribution des Textsatzsystems TEX
			\item TeXStudio - Entwicklungsumgebung für \LaTeX
			\item Balsamiq - Programm zur Erstellung von Mockups
			\item microTool inStep - Projektverwaltungstool für Arbeitsmaschine 2
		\end{itemize}
	}
	\itemd{Personal}{
		\begin{itemize}
			\item Projektbetreuer zur Unterstützung
		\end{itemize}
	}
\end{itemize}

\subsection{Entwicklungsprozess}
\label{sec:Entwicklungsprozess}

Die ausgewählte Vorgehensweise ist das Wasserfall-Modell\footnote{Wasserfall-Modell nach W. Royce}.

\begin{enumerate}
	\itemd{Systemanforderungen:}{ Alle Anforderungen, die selbst nicht direkt das Software-Produkt betreffen, werden zunächst festgelegt. Dazu zählen:}
		\begin{itemize}
			\item Preis
			\item Verfügbarkeit
			\item Sicherheitsaspekte
			\item Dokumentation
		\end{itemize}
	
	\itemd{Softwareanforderungen:}{Alle Anforderung an die Software selber werden definiert. Jegliche Funktionen, Interaktionen und Besonderheiten werden konkretisiert, so dass sich aus den Systemanforderungen und Softwareanforderungen das Lastenheft ergibt.}
	
	\itemd{Analyse:}{Anforderungen aus Lastenheft und Ist-Zustand der Situation werden analysiert, so dass diese in ein Pflichtenheft umformuliert werden können. Die Wirtschaftlichkeit eines Projektes wird ebenfalls hier geprüft.}

	\itemd{Entwurf:}{Es wird das Datenmodell, die Architektur und die Schnittstellen zu anderen Anwendungen herausgearbeitet. Zwischenergebnisse können sein:}
		\begin{itemize}
			\item \ac{ERM}
			\item \acs{UML}-Diagramme (Klassendiagramm, Sequenzdiagramm, Anwendungsfalldiagramm \usw)
			\item Mockups zur \ac{GUI}
			\item Schnittstellen-Verzeichnis
		\end{itemize}
	
	\itemd{Implementierung:}{Umsetzung der Funktionalitäten nach Pflichtenheft und Entwurf in eine lauffähige Anwendung.}
	
	\itemd{Test/Qualitätssicherung:}{Es wird nach der Implementierungsphase die Software auf Fehler, Schwachstellen und Unstimmigkeiten überprüft. weitverbreitete Testmethoden:}
		\begin{itemize}
			\item Komponententests (Unit-Test)
			\item Modultests
			\item Systemtests
			\item Integrationstests
		\end{itemize}
	
	\itemd{Inbetriebnahme:}{letzter Schritt in der Softwareentwicklung. Bei fehlerloser Überprüfung kann die Anwendung abgenommen werden und in Produktion gehen.}
\end{enumerate} 

Alle Schritte im Wasserfall-Modell sind sequentiell zu bearbeiten, \dahe Kein Schritt darf vorgezogen werden oder übersprungen werden. Es dürfen nur Phasen zurückgesprungen werden, Wenn eine Phase nicht abgeschlossenen werden kann.

%ProjektPhasen
%\begin{itemize}
%	\item In welchem Zeitraum und unter welchen Rahmenbedingungen (\zB Tagesarbeitszeit) findet das Projekt statt?
%	\item Verfeinerung der Zeitplanung, die bereits im Projektantrag vorgestellt wurde.
%\end{itemize}

%ProjektAbweichungen
%\begin{itemize}
%	\item Sollte es Abweichungen zum Projektantrag geben (\zB Zeitplanung, Inhalt des Projekts, neue Anforderungen), müssen diese explizit aufgeführt und begründet werden.
%\end{itemize}

%RessourcenPlanung
%\begin{itemize}
%	\item Detaillierte Planung der benötigten Ressourcen (Hard-/Software, Räumlichkeiten \usw).
%	\item \Ggfs sind auch personelle Ressourcen einzuplanen (\zB unterstützende Mitarbeiter).
%	\item Hinweis: Häufig werden hier Ressourcen vergessen, die als selbstverständlich angesehen werden (\zB PC, Büro). 
%\end{itemize}

%Entwicklungsprozess
%\begin{itemize}
%	\item Welcher Entwicklungsprozess wird bei der Bearbeitung des Projekts verfolgt (\zB Wasserfall, agiler Prozess)?
%\end{itemize}
% !TEX root = ../Projektdokumentation.tex

\newpage
\section{Analysephase} 
\label{sec:Analysephase}


\subsection{Ist-Analyse} 
\label{sec:IstAnalyse}

Die Betriebsprüfer müssen die Dokumente einzeln über den Dokumenttypenbaum aktualisieren. Dies kann sehr aufwendig sein, da manche Berechnungen in mehreren Dokumenten aufgelistet werden und sind diese erst ausfindig zu machen, um dann aktualisiert werden zu können.

Diese Funktionalität ist von den Betriebsprüfern dringend erwünscht, da es eine enorme Zeitersparnis für sie ergeben würde.
\subsection{Wirtschaftlichkeitsanalyse}
\label{sec:Wirtschaftlichkeitsanalyse}

\subsubsection{\gqq{Make or Buy}-Entscheidung}
\label{sec:MakeOrBuyEntscheidung}

Da die Entwicklung der \acs{CBP-AD} ein internes Projekt der \acs{DRV} ist und nur eine Funktionserweiterung ist, lässt sich kein fertiges Produkt finden, dass alle Anforderungen, vor allem fachliche, erfüllt. Daher wird dieses Projekt in Eigenentwicklung von der IT-Abteilung der \acs{DRV} umgesetzt.

\subsubsection{Projektkosten}
\label{sec:Projektkosten}

Die Kosten für die Durchführung des Projekts setzen sich sowohl aus Personal-, als auch aus Ressourcenkosten zusammen.
Laut Tarifvertrag\footnote{\href{http://www.oeffentlichen-dienst.de/auszubildende/25-tarifvertrag-fuer-den-auszubildende/90-tvaoed-besonderer-teil-bbig.html}{Link zum Tarifvertrag}} verdient ein Auszubildender im dritten Lehrjahr pro Monat \eur{949,02} Brutto. 

\begin{eqnarray}
8 \mbox{ h/Tag} \cdot 220 \mbox{ Tage/Jahr} = 1760 \mbox{ h/Jahr}\\
\eur{949.02}\mbox{/Monat} \cdot 12,9 \mbox{ Monate/Jahr} \approx \eur{12242,36} \mbox{/Jahr}\\
\frac{\eur{12242,36} \mbox{/Jahr}}{1760 \mbox{ h/Jahr}} \approx \eur{6,96}\mbox{/h}
\end{eqnarray}

Es ergibt sich also ein Stundenlohn von \eur{6,96}. 
Die Durchführungszeit des Projekts beträgt 70 Stunden. Für die Nutzung von Ressourcen\footnote{Räumlichkeiten, Arbeitsplatzrechner etc.} wird 
ein pauschaler Stundensatz von \eur{15} angenommen. Für die anderen Mitarbeiter wird pauschal ein Stundenlohn von \eur{25} angenommen. 
Eine Aufstellung der Kosten befindet sich in Tabelle~\ref{tab:Kostenaufstellung} und sie betragen insgesamt \eur{1697,20}.
\tabelle{Kostenaufstellung}{tab:Kostenaufstellung}{Kostenaufstellung.tex}


\subsubsection{Amortisationsdauer}
\label{sec:Amortisationsdauer}

Bei einer Zeiteinsparung von 10 Minuten am Tag für jeden der 4000 Betriebsprüfer und 220 Arbeitstagen im Jahr ergibt sich eine gesamte Zeiteinsparung von 
\begin{eqnarray}
4000 \cdot 220 \mbox{ Tage/Jahr} \cdot 10 \mbox{ min/Tag} = 8800000 \mbox{ min/Jahr} \approx 146667 \mbox{ h/Jahr} 
\end{eqnarray}

Dadurch ergibt sich eine jährliche Einsparung von 
\begin{eqnarray}
146667 \mbox{h} \cdot \eur{(25 + 15)}{\mbox{/h}} = \eur{5866680}
\end{eqnarray}

Die Amortisationszeit beträgt also $\frac{\eur{1697,20}}{\eur{5866680}\mbox{/Jahr}} \approx 0,0002 \mbox{ Jahre} \approx 2 \mbox{ Stunden}$.


%UseCases - entfällt wegen nur einem


\subsection{Lastenheft/Fachkonzept}
\label{sec:Lastenheft}

\paragraph{Beispiel}
Ein Beispiel für ein Lastenheft findet sich im \Anhang{app:Lastenheft}. 

%\Zwischenstand{Analysephase}{Analyse}

%Analyse
%\begin{itemize}
%	\item Wie ist die bisherige Situation (\zB bestehende Programme, Wünsche der Mitarbeiter)?
%	\item Was gilt es zu erstellen/verbessern?
%\end{itemize}

%WirtschaftlichkeitsAnalyse
%\begin{itemize}
%	\item Lohnt sich das Projekt für das Unternehmen?
%\end{itemize}

%%Make-Or-Buy-Entscheidung
%\begin{itemize}
%	\item Gibt es vielleicht schon ein fertiges Produkt, dass alle Anforderungen des Projekts abdeckt?
%	\item Wenn ja, wieso wird das Projekt trotzdem umgesetzt?
%\end{itemize}

%%ProjektKosten
%\begin{itemize}
%	\item Welche Kosten fallen bei der Umsetzung des Projekts im Detail an (\zB Entwicklung, Einführung/Schulung, Wartung)?
%\end{itemize}

%%Amortisationsdauer
%\begin{itemize}
%	\item Welche monetären Vorteile bietet das Projekt (\zB Einsparung von Lizenzkosten, Arbeitszeitersparnis, bessere Usability, Korrektheit)?
%	\item Wann hat sich das Projekt amortisiert?
%\end{itemize}

%Nutzwertanalyse
%\begin{itemize}
%	\item Darstellung des nicht-monetären Nutzens (\zB Vorher-/Nachher-Vergleich anhand eines Wirtschaftlichkeitskoeffizienten). 
%\end{itemize}

%Anwendungsfälle
%\begin{itemize}
%	\item Welche Anwendungsfälle soll das Projekt abdecken?
%	\item Einer oder mehrere interessante (!) Anwendungsfälle könnten exemplarisch durch ein Aktivitätsdiagramm oder eine \ac{EPK} detailliert beschrieben werden. 
%\end{itemize}

%Qualitätsanforderungen
%\begin{itemize}
%	\item Welche Qualitätsanforderungen werden an die Anwendung gestellt (\zB hinsichtlich Performance, Usability, Effizienz \etc (siehe \citet{ISO9126}))?
%\end{itemize}

%Lastenheft/Fachkonzept
%\begin{itemize}
%	\item Auszüge aus dem Lastenheft/Fachkonzept, wenn es im Rahmen des Projekts erstellt wurde.
%	\item Mögliche Inhalte: Funktionen des Programms (Muss/Soll/Wunsch), User Stories, Benutzerrollen
%\end{itemize}

%Zwischenstand

% !TEX root = ../Projektdokumentation.tex
\section{Entwurfsphase} 
\label{sec:Entwurfsphase}

\subsection{Zielplattform}
\label{sec:Zielplattform}

\begin{itemize}
	\item Beschreibung der Kriterien zur Auswahl der Zielplattform (\ua Programmiersprache, Datenbank, Client/Server, Hardware).
\end{itemize}


\subsection{Architekturdesign}
\label{sec:Architekturdesign}
\begin{itemize}
	\item Beschreibung und Begründung der gewählten Anwendungsarchitektur (\zB \acs{MVC}).
	\item \Ggfs Bewertung und Auswahl von verwendeten Frameworks sowie \ggfs eine kurze Einführung in die Funktionsweise des verwendeten Frameworks.
\end{itemize}

\paragraph{Beispiel}
Anhand der Entscheidungsmatrix in Tabelle~\ref{tab:Entscheidungsmatrix} wurde für die Implementierung der Anwendung das \acs{PHP}-Framework Symfony\footnote{\Vgl \citet{Symfony}.} ausgewählt. 

\tabelle{Entscheidungsmatrix}{tab:Entscheidungsmatrix}{Nutzwert.tex}


\subsection{Entwurf der Benutzeroberfläche}
\label{sec:Benutzeroberflaeche} 
\begin{itemize}
	\item Entscheidung für die gewählte Benutzeroberfläche (\zB GUI, Webinterface).
	\item Beschreibung des visuellen Entwurfs der konkreten Oberfläche (\zB Mockups, Menüführung).
	\item \Ggfs Erläuterung von angewendeten Richtlinien zur Usability und Verweis auf Corporate Design.
\end{itemize}

\paragraph{Beispiel}
Beispielentwürfe finden sich im \Anhang{app:Entwuerfe}.


\subsection{Datenmodell}
\label{sec:Datenmodell}

\begin{itemize}
	\item Entwurf/Beschreibung der Datenstrukturen (\zB \acs{ERM} und/oder Tabellenmodell, \acs{XML}-Schemas) mit kurzer Beschreibung der wichtigsten (!) verwendeten Entitäten.
\end{itemize}

\paragraph{Beispiel}
In \Abbildung{ER} wird ein \ac{ERM} dargestellt, welches lediglich Entitäten, Relationen und die dazugehörigen Kardinalitäten enthält. 

\begin{figure}[htb]
\centering
\includegraphicsKeepAspectRatio{ERDiagramm.pdf}{0.6}
\caption{Vereinfachtes ER-Modell}
\label{fig:ER}
\end{figure} 


\subsection{Geschäftslogik}
\label{sec:Geschaeftslogik}

\begin{itemize}
	\item Modellierung und Beschreibung der wichtigsten (!) Bereiche der Geschäftslogik (\zB mit Kom\-po\-nen\-ten-, Klassen-, Sequenz-, Datenflussdiagramm, Programmablaufplan, Struktogramm, \ac{EPK}).
	\item Wie wird die erstellte Anwendung in den Arbeitsfluss des Unternehmens integriert?
\end{itemize}

\paragraph{Beispiel}
Ein Klassendiagramm, welches die Klassen der Anwendung und deren Beziehungen untereinander darstellt kann im \Anhang{app:Klassendiagramm} eingesehen werden.

\Abbildung{Modulimport} zeigt den grundsätzlichen Programmablauf beim Einlesen eines Moduls als \ac{EPK}.
\begin{figure}[htb]
\centering
\includegraphicsKeepAspectRatio{modulimport.pdf}{0.9}
\caption{Prozess des Einlesens eines Moduls}
\label{fig:Modulimport}
\end{figure}


\subsection{Maßnahmen zur Qualitätssicherung}
\label{sec:Qualitaetssicherung}
\begin{itemize}
	\item Welche Maßnahmen werden ergriffen, um die Qualität des Projektergebnisses (siehe Kapitel~\ref{sec:Qualitaetsanforderungen}: \nameref{sec:Qualitaetsanforderungen}) zu sichern (\zB automatische Tests, Anwendertests)?
	\item \Ggfs Definition von Testfällen und deren Durchführung (durch Programme/Benutzer).
\end{itemize}


\subsection{Pflichtenheft/Datenverarbeitungskonzept}
\label{sec:Pflichtenheft}
\begin{itemize}
	\item Auszüge aus dem Pflichtenheft/Datenverarbeitungskonzept, wenn es im Rahmen des Projekts erstellt wurde.
\end{itemize}

\paragraph{Beispiel}
Ein Beispiel für das auf dem Lastenheft (siehe Kapitel~\ref{sec:Lastenheft}: \nameref{sec:Lastenheft}) aufbauende Pflichtenheft ist im \Anhang{app:Pflichtenheft} zu finden.


\Zwischenstand{Entwurfsphase}{Entwurf}

% !TEX root = ../Projektdokumentation.tex
\section{Implementierungsphase} 
\label{sec:Implementierungsphase}

\subsection{Implementierung der Geschäftslogik}
\label{sec:ImplementierungGeschaeftslogik}
%
%\begin{itemize}
%	\item Beschreibung des Vorgehens bei der Umsetzung/Programmierung der entworfenen Anwendung.
%	\item \Ggfs interessante Funktionen/Algorithmen im Detail vorstellen, verwendete Entwurfsmuster zeigen.
%	\item Quelltextbeispiele zeigen.
%	\item Hinweis: Wie in Kapitel~\ref{sec:Einleitung}: \nameref{sec:Einleitung} zitiert, wird nicht ein lauffähiges Programm bewertet, sondern die Projektdurchführung. Dennoch würde ich immer Quelltextausschnitte zeigen, da sonst Zweifel an der tatsächlichen Leistung des Prüflings aufkommen können.
%\end{itemize}
%
%\paragraph{Beispiel}
%Die Klasse \texttt{Com\-par\-ed\-Na\-tu\-ral\-Mo\-dule\-In\-for\-ma\-tion} findet sich im \Anhang{app:CNMI}.  


\Zwischenstand{Implementierungsphase}{Implementierung}

%Implementierung der Datenstrukturen

%\begin{itemize}
%	\item Beschreibung der angelegten Datenbank (\zB Generierung von \acs{SQL} aus Modellierungswerkzeug oder händisches Anlegen), \acs{XML}-Schemas \usw.
%\end{itemize}


%Implementierung der Benutzeroberfläche

%\begin{itemize}
%	\item Beschreibung der Implementierung der Benutzeroberfläche, falls dies separat zur Implementierung der Geschäftslogik erfolgt (\zB bei \acs{HTML}-Oberflächen und Stylesheets).
%	\item \Ggfs Beschreibung des Corporate Designs und dessen Umsetzung in der Anwendung.
%	\item Screenshots der Anwendung
%\end{itemize}


%Implementierung der Geschäftslogik

%\begin{itemize}
%	\item Beschreibung des Vorgehens bei der Umsetzung/Programmierung der entworfenen Anwendung.
%	\item \Ggfs interessante Funktionen/Algorithmen im Detail vorstellen, verwendete Entwurfsmuster zeigen.
%	\item Quelltextbeispiele zeigen.
%	\item Hinweis: Wie in Kapitel~\ref{sec:Einleitung}: \nameref{sec:Einleitung} zitiert, wird nicht ein lauffähiges Programm bewertet, sondern die Projektdurchführung. Dennoch würde ich immer Quelltextausschnitte zeigen, da sonst Zweifel an der tatsächlichen Leistung des Prüflings aufkommen können.
%\end{itemize}
% !TEX root = ../Projektdokumentation.tex
\section{Abnahmephase} 
\label{sec:Abnahmephase}

\subsection{Komponententest}

Der in \ref{sec:Qualitaetssicherung} beschriebene Komponententest befindet sich im \Anhang{app:Test}.
Die Ausführung dieses Tests führt zu folgendem Resultat.

\begin{figure}[h]
	\centering
	\includegraphicsKeepAspectRatio{UnitTestSession.png}{1}
	\caption{Ausführung des Komponententests}
\end{figure}

\subsection{Abnahme der IT-Abteilung}

Bei erfolgreichen Komponententests prüft ein zweiter Mitarbeiter mittels Code-Review, ob auffällige Stellen, Flüchtigkeitsfehler oder ähnliches vorliegen. Bei vorläufiger Abnahme ist die Erweiterung soweit zum nächsten sogenannten Verbundtest\footnote{System- und Benutzeroberflächentests aller Programme durch die Fachabteilung} in die Anwendung integriert zu werden. Dann erst kann komplett bestätigt werden, ob sich alle Funktionen der Anwendung weiterhin fehlerfrei ausführen lassen und sich an die Richtlinien zur Barrierefreiheit und Usability gehalten wird.

\Zwischenstand{Abnahmephase}{Abnahme}

%% !TEX root = ../Projektdokumentation.tex
\section{Einführungsphase}
\label{sec:Einfuehrungsphase}

\begin{itemize}
	\item Welche Schritte waren zum Deployment der Anwendung nötig und wie wurden sie durchgeführt (automatisiert/manuell)?
	\item Wurden \ggfs Altdaten migriert und wenn ja, wie?
	\item Wurden Benutzerschulungen durchgeführt und wenn ja, Wie wurden sie vorbereitet?
\end{itemize}


\Zwischenstand{Einführungsphase}{Einfuehrung}

% !TEX root = ../Projektdokumentation.tex
\section{Dokumentation}
\label{sec:Dokumentation}

\begin{itemize}
	\item Wie wurde die Anwendung für die Benutzer/Administratoren/Entwickler dokumentiert (\zB Benutzerhandbuch, \acs{API}-Dokumentation)?
	\item Hinweis: Je nach Zielgruppe gelten bestimmte Anforderungen für die Dokumentation (\zB keine IT-Fachbegriffe in einer Anwenderdokumentation verwenden, aber auf jeden Fall in einer Dokumentation für den IT-Bereich).
\end{itemize}

\paragraph{Beispiel}
Ein Ausschnitt aus der erstellten Benutzerdokumentation befindet sich im \Anhang{app:BenutzerDoku}.
Die Entwicklerdokumentation wurde mittels PHPDoc\footnote{Vgl. \cite{phpDoc}} automatisch generiert. Ein beispielhafter Auszug aus der Dokumentation einer Klasse findet sich im \Anhang{app:Doc}. 

\Zwischenstand{Dokumentation}{Dokumentation}

% !TEX root = ../Projektdokumentation.tex
\section{Fazit} 
\label{sec:Fazit}

\subsection{Soll-/Ist-Vergleich}
\label{sec:SollIstVergleich}

Wie in Tabelle~\ref{tab:Vergleich} zu erkennen ist, konnte die Zeitplanung bis auf wenige Ausnahmen eingehalten werden. Die Analysephase brauchte weniger Zeit als geplant, weil die IT-Abteilung schon bereits den größten Teil der fachlichen Analyse vorbereitet hatte. Das Erstellen der Dokumentation hingegen brauchte etwas mehr Zeit als geplant, da die Verwendung von LateX zu einigen Komplikation geführt hatte.
\tabelle{Soll-/Ist-Vergleich}{tab:Vergleich}{Zeitnachher.tex}


\subsection{Lessons Learned}
\label{sec:LessonsLearned}

Durch Projekte wie dieses wird einem erst bewusst, wie wichtig Anforderungen und Spezifikationen sind, da sie zum Teil die einzigen Anhaltspunkte zur Feststellung des Entwicklungsfortschritts sind.

\subsection{Ausblick}
\label{sec:Ausblick}

Nach erfolgreichem Abschließen der Verbundtests können die Anwendungen in Produktion gehen.

Es sind keine Erweiterungen der Funktionalität geplant.


% Literatur ------------------------------------------------------------------
\clearpage
\renewcommand{\refname}{Literaturverzeichnis}
\bibliography{Bibliographie}
\bibliographystyle{Allgemein/natdin} % DIN-Stil des Literaturverzeichnisses
% !TEX root = Projektdokumentation.tex
\clearpage
\addsec{Eidesstattliche Erklärung}

% Hinweis: die eidesstattliche Erklärung ist ggfs. an die Richtlinie der IHK anzupassen

Ich, \autorName, versichere hiermit, dass ich meine \textbf{\betreff} mit dem
Thema
\begin{quote}
\textit{\kompletterTitel}
\end{quote}
selbständig verfasst und keine anderen als die angegebenen Quellen und Hilfsmittel benutzt habe,
wobei ich alle wörtlichen und sinngemäßen Zitate als solche gekennzeichnet habe. Die Arbeit
wurde bisher keiner anderen Prüfungsbehörde vorgelegt und auch nicht veröffentlicht.\\[6ex]

\abgabeOrt, den \abgabeTermin


\rule[-0.2cm]{5.5cm}{0.5pt}

\textsc{\autorName}


% Anhang ---------------------------------------------------------------------
\clearpage
\appendix
\pagenumbering{roman}
% !TEX root = Projektdokumentation.tex
\section{Anhang}

\subsection{Detaillierte Zeitplanung}
\label{app:Zeitplanung}
\tabelleAnhang{ZeitplanungKomplett}

\clearpage
\subsection{Klassendiagramm}
\label{app:Klassendiagramm}

Die grün umrandeten Klassen sind verwendete, bereits existierende Klassen der Code-Basis.
Alle schwarz umrandeten Klassen sind durch dieses Projekt entstanden.

\begin{figure}[htb]
\centering
\includegraphicsKeepAspectRatio{Klassendiagramm.png}{1.2}
\caption{Vollständiges Klassendiagramm}
\end{figure}

\subsection{Sequenzdiagramm}
\label{app:Sequenzdiagramm}

\begin{figure}[htb]
	\centering
	\includegraphicsKeepAspectRatio{Sequenzdiagramm.png}{1}
	\caption{Sequenzdiagramm}
\end{figure}

\clearpage
\subsection{Screenshots der Anwendung}
\label{app:Screenshots}

\begin{figure}[htb]
	\centering
	\includegraphicsKeepAspectRatio{Dokumentenbaum.png}{0.7}
	\caption{Dokumentenbaum mit den verschiedenen Dokumenttypen sortiert nach Prüfgebieten}
\end{figure}

\begin{figure}[htb]
	\centering
	\includegraphicsKeepAspectRatio{Menueeintrag.png}{0.7}
	\caption{Menüeintrag zum Anstoß des Globalen Aktualisierens}
\end{figure}


\clearpage
\subsection{Testfälle}
\label{app:Test}

\begin{figure}[htb]
	\centering
	\includegraphicsKeepAspectRatio{UnitTest.png}{1}
	\caption{Komponententest}
\end{figure}

%\lstinputlisting[language=php]{Listings/tests.php}
%\clearpage
%\begin{figure}[htb]
%\centering
%\includegraphicsKeepAspectRatio{testcase.jpg}{1}
%\caption{Aufruf des Testfalls auf der Konsole}
%\end{figure}



\end{document}